\section*{Bilan du maintien en condition opérationnelle}
Depuis la mise en production du produit minimum viable le 10~août~2026 (\texttt{v0.6.0}), l'application est exploitée en ligne et fait l'objet d'un suivi opérationnel. Le dispositif de maintien a été mis en place \emph{dès} la mise en service~: supervision et mesure d'audience (Rybbit) couplées à un endpoint de santé, veille et mise à jour des dépendances (Renovate), analyse continue de qualité et de sécurité (SonarQube, Semgrep, Trivy). Ce dispositif a démontré sa valeur sur le premier incident~: l'anomalie \texttt{\#142}, détectée par la supervision, a été consignée puis corrigée et déployée via la chaîne CI/CD (\texttt{v0.6.1}), la disponibilité étant restée maîtrisée sur la période.

\section*{Points forts}
Trois éléments ressortent. La \textbf{réutilisation intégrale de la chaîne CI/CD} du Bloc~2 permet de livrer un correctif rapidement et sans régression, sans dispositif spécifique au maintien. La \textbf{supervision est légère mais ciblée}~: en surveillant la santé applicative réelle (endpoint \texttt{/health}) plutôt qu'une simple page, elle évite les faux négatifs tout en restant proportionnée à un déploiement mono-serveur. Enfin, la \textbf{traçabilité est complète}, de la détection à la résolution jusqu'au journal des versions.

\section*{Points d'amélioration et perspectives}
Le recul d'exploitation reste court~; le dispositif gagnera à être consolidé dans la durée. Plusieurs chantiers sont identifiés~: industrialiser la supervision (métriques centralisées et tableaux de bord au-delà de l'uptime), automatiser davantage l'évaluation des mises à jour majeures, et formaliser le canal de support. Les axes d'amélioration argumentés au chapitre~3 — allègement du frontend, résorption de la dette du moteur de tournoi, notifications de résultats — ainsi que la conformité RGPD constituent la feuille de route des prochaines versions.

\section*{Retour d'expérience personnel}
Ce bloc m'a fait passer de la posture de développeur à celle de mainteneur~: penser l'application non plus seulement à travers sa réalisation, mais à travers son exploitation dans le temps, et arbitrer en permanence entre correctifs, évolutions et dette. J'ai mesuré qu'une part du maintien consiste autant à surveiller, expliquer et tracer qu'à coder. Je remercie Roméo Chevrier pour son accompagnement sur la mise en production, Farès Sioni pour son aide sur la structure du dossier, Paul-Lucas Estrada pour son implication en tant que commanditaire, et mon tuteur Fabrice Gascoin pour sa relecture et ses conseils.
